\chapter{Вывод}
\label{ch:chap3}

В ходе проделанной работы я познакомился c идеей кластеризации и построил модель на основе методу k-means и проверил ее эффективность.
% \section{Контрольные вопросы}

% \begin{enumerate} 
%     \item Какие инструментальные средства используются для организации рабочего места специалиста Data Science? 
%     \item Какие библиотеки Python используются для работы в области машинного обучения? Дайте краткую характеристику каждой библиотеке.
%     \item Почему при реализации систем машинного обучения широкое распространение получили библиотеки Python?
% \end{enumerate}

% Ответы:

% \begin{enumerate}
%     \item Для организации рабочего места специалиста Data Science используется: язык программирования, IDE, библиотеки и фреймворки, базы данных. 
%     \item В области машинного обучения для Python используются библиотеки numpy, pandas, matplotlib (для первичного анализа данных).
%     \item Python популярен в машинном обучении за счет своего удобства, читаемости кода, большого кол-ва библиотек и фреймворков, кроссплатформенности. Многие модули Python на самом деле написаны на C/C++, что делает код очень быстрым.
% \end{enumerate}

\endinput